\justifying
\NSFCBodyText

\subsubsection{研究方法}

以完整运输任务 $i$ 和完整仓储事件 $j$ 为不可拆分数据单元，$t$ 为决策时刻，$m$ 为运输模态，$n$ 为仓储信息源。运输观测 $X_t^{(m)}$ 包括振动、倾斜、电子封志、光感和包装图像；$M_t^{(m)}$、$\Delta_t^{(m)}$ 与 $q_t^{(m)}$ 分别表示可用性、时间偏移和观测质量。首先分别提取模态内特征并估计时变可靠性：
\[
\begin{aligned}
h_t^{(m)}&=f_m\!\left(X_{t-L_m+1:t}^{(m)},\Delta_t^{(m)},q_t^{(m)}\right),\\
r_t^{(m)}&=\sigma\!\left(g_m(h_t^{(m)},q_t^{(m)},M_t^{(m)})\right),\quad 0\le r_t^{(m)}\le1.
\end{aligned}
\tag{1}
\]
$h_t^{(m)}$ 为时序或视觉表征，$r_t^{(m)}$ 随噪声、缺失和跨模态冲突变化。通过受控信号噪声、时间错位、图像遮挡与模态删除，检验可靠性是否随扰动单调下降。

可用模态经可靠性加权形成统一货物状态，并输出事件类别、包装损伤位置与不确定性：
\[
z_t=\frac{\sum_m M_t^{(m)}r_t^{(m)}h_t^{(m)}}
{\sum_m M_t^{(m)}r_t^{(m)}+\varepsilon},\qquad
(\hat y_t^{c},\hat b_t,\hat u_t)=F_\theta(z_{1:t}).
\tag{2}
\]
其中 $\hat y_t^{c}$ 表示冲击/倾斜、非法开启等事件概率，$\hat b_t$ 表示包装破损或变形位置，$\hat u_t$ 表示预测不确定性。式（2）分别对应事件 AUPRC/F1、图像 mAP/IoU 和 Brier/ECE。

货物感知模型的联合优化目标为
\[
\mathcal L_{1}=\mathcal L_{\mathrm{evt}}
+\lambda_{\mathrm{img}}\mathcal L_{\mathrm{damage}}
+\lambda_{\mathrm{align}}\mathcal L_{\mathrm{align}}
+\lambda_{\mathrm{cal}}\mathcal L_{\mathrm{cal}}.
\tag{3}
\]
四项分别约束事件识别、包装损伤定位、异步模态对齐和概率校准，各权重只在开发集确定。若完整运输任务外层测试中联合模型不能同时改善核心任务指标与校准，或相同容量的无可靠性模型获得同等结果，则 H1 失败。

对于弱网边缘协同，端侧状态定义为 $s_t=(\hat y_t^{c},\hat u_t,N_t,Q_t,B_t,C_t,A_t^{\mathrm{rx}})$，分别包含风险、网络、缓存队列、电池、算力和接收端确认的断点状态；动作包括采样粒度、本地推理深度、压缩级别、缓存优先级、上传/续传和工作周期。策略在可靠送达与尾时延约束下联合优化：
\[
\begin{aligned}
\min_{\pi}\;&\mathbb E_{\pi}\sum_t
\left(\lambda_m C_t^{\mathrm{miss}}+\lambda_d T_t
+\lambda_b Z_t^{\mathrm{byte}}+\lambda_e E_t\right),\quad a_t=\pi(s_t),\\
\mathrm{s.t.}\;&\Pr(D_t=1\mid R_t\ge r_H)\ge1-\epsilon,\quad
Q_{0.95}(T_t^{\mathrm{alarm}})\le\tau .
\end{aligned}
\tag{4}
\]
$C_t^{\mathrm{miss}}$ 为高风险漏报代价，$T_t$、$Z_t^{\mathrm{byte}}$、$E_t$ 和 $D_t$ 分别为告警时延、实际传输字节、端侧能耗和按时送达指示。式（4）把低功耗与可靠上传置于同一目标；业务阈值无事实依据时，采用相对冻结基线的不劣口径，不预设数值。

候选证据 $e$ 的风险信息价值和动作判据写为
\[
\begin{aligned}
V_t(e,a)=&\;p_t^{\mathrm{on}}(e,a)e^{-\kappa\mathrm{AoI}_t(e)}
\left[\mathcal R(b_t^r)-\mathbb E\mathcal R\!\left(\mathrm{Update}(b_t^r,e)\right)\right]\\
&-\eta_b Z_t^{\mathrm{byte}}(e,a)-\eta_e E_t(e,a)-\eta_c C_t(e,a),\\
a_t^*&=&\arg\max_{a\in\mathcal A}V_t(e,a).
\end{aligned}
\tag{5}
\]
$b_t^r$ 为接收端已确认信息形成的风险信念，$\mathcal R$ 为告警决策风险，$p_t^{\mathrm{on}}$ 为按时送达概率，$C_t(e,a)$ 为计算开销；$\mathcal A$ 包含端侧摘要、缓存、发送和续传。缓存保存优先级、校验值与已确认偏移，重连后从断点续传；高风险证据抢占队列，稳定低价值片段在端侧摘要。若同一事件流和网络轨迹上不能在相同安全约束下减少字节/能耗或提高高风险送达率，则 H2 失败。

仓储观测 $W_t^{(n)}$ 包括温度、湿度、安防视频、烟雾和门禁。针对其时标差异建立多源关系图，风险状态更新为
\[
H_t^{w}=G_\phi\!\left(H_{t-1}^{w},
\left\{r_t^{(n)}f_n(W_{t-L_n+1:t}^{(n)})\right\}_{n=1}^{K},
\rho_t^{\mathrm{tr}}\right),\qquad
\hat p_t^{w}=\operatorname{softmax}(U H_t^{w}).
\tag{6}
\]
$H_t^{w}$ 为仓储风险状态，$\hat p_t^{w}$ 为正常、环境异常、安防异常和复合风险概率；$\rho_t^{\mathrm{tr}}$ 为可选运输风险先验，无货物级配对记录时置为空，仓储模型仍以四类信息独立训练和验证。

仓储模型同时约束事件识别、提前预警、概率校准和跨源时序关系：
\[
\mathcal L_{3}=\mathcal L_{\mathrm{event}}
+\mu_{\mathrm{early}}\mathcal L_{\mathrm{early}}
+\mu_{\mathrm{cal}}\mathcal L_{\mathrm{cal}}
+\mu_{\mathrm{rel}}\mathcal L_{\mathrm{relation}}.
\tag{7}
\]
$\mathcal L_{\mathrm{early}}$ 对风险发生前的正确提示进行约束，$\mathcal L_{\mathrm{relation}}$ 描述温湿度缓变与烟雾、门禁、视频突发的共现和先后关系。通过来源删除、时间置乱、固定可靠性和无运输先验消融检验关系增益。

校准风险进一步映射为分级预警与响应动作：
\[
k_t^*=\min\left\{k:\Pr(Y_t^{w}\ge k\mid H_t^{w})\ge\tau_k\right\},\qquad
u_t^w=\Gamma(k_t^*,\mathcal E_t,d_t).
\tag{8}
\]
$k_t^*$ 为风险等级，$\mathcal E_t$ 为主要证据来源，$d_t$ 为持续时间，$\Gamma$ 输出记录、复核、联动通知或应急处置建议。阈值 $\tau_k$ 在校准集冻结；若完整仓储事件外层留出中不优于独立阈值、单源或无关系模型，或时间/来源置乱不降低收益，则 H3 失败。

\subsubsection{技术路线}

技术路线依次为：制定运输任务和仓储事件数据字典；采集/构造五类货物观测、弱网资源状态和四类仓储信息；完成时间对齐、质量估计和可靠性加权表征；识别运输异常和包装损伤；计算风险信息价值并联合执行边缘推理、压缩、缓存、断点续传、上传和动态工作周期；建立仓储多源风险演化与分级预警；最后在完整任务/事件外层测试中开展三条主线的独立反证与端到端联调。内容三不以运输先验为前置条件，内容一和内容二也分别设置独立基线，避免联合系统共用测试集调参。

\subsubsection{实验手段}

货物感知实验通过可控振动/冲击、倾斜、封志开启、光感变化和包装破损/变形获得事件真值，与单模态、静态早/晚融合、缺失填补和普通 mask-aware 模型比较；公开工业异常数据仅用于视觉预训练或方法对照。弱网实验记录网络可用性、带宽、丢包、时延、ACK、缓存偏移、重传、推理时延和能耗，在相同事件流、网络轨迹和初始电量下比较周期上传、阈值触发、仅信息年龄及固定压缩/工作周期。仓储实验在可控环境中组织温湿度缓变、烟雾、异常门禁、视频安防及其组合事件，与独立阈值、单源分类和无时序关系模型比较。

所有数据记录测量时点、可用时点、来源和完整任务/事件标识，只使用可用时点不晚于预测时点的信息；按运输任务或仓储事件划分训练、开发、校准和封存外层测试，禁止相邻窗口跨集合。主要结果报告事件 AUPRC/F1、mAP/IoU、Brier/ECE、上传成功率、高风险漏报、P95 告警时延、传输字节、单位任务能耗、仓储误报率和预警提前量，并给出按完整任务/事件聚类的置信区间。

\subsubsection{关键技术}

关键技术包括：一是面向异步和不完备观测的模态可靠性动态估计，使五类货物证据在不同扰动下具有可解释权重；二是以风险信息价值为纽带联合边缘推理、压缩、缓存、断点续传、可靠上传和动态工作周期，使低功耗管理不以牺牲高风险事件为代价；三是面向仓储缓变与突发信号的不同时标关系建模和概率校准，使预警等级能够追溯到证据来源与持续时间；四是以完整任务/事件划分、可用时点约束、同轨迹配对和置乱负对照防止数据泄漏与机制误判。

\subsubsection{可行性分析}

申请人已有多模态信息处理、视觉理解、时间序列建模和深度学习研究积累，可支撑内容一和内容三的表征、检测与动态建模；既有工程研发经历可支撑端侧算法实现和实验链路组织。项目采用“受控实验建立内部有效性—弱网轨迹和固定设备验证资源机制—条件具备后开展现场外部验证”的分层路线，不以尚未确认的合作、设备或真实样本作为主假设成立的唯一条件。

主要风险及应对为：货损标签不足时，以受控事件和弱监督正常表征获得真值，并限定外推范围；现场 2G/3G 轨迹不足时，公开参数并进行多强度可复现回放；资源协同无联合收益时，保留冻结安全基线并报告 H2 失败；仓储复合事件不足时，先完成单事件和可控组合事件分层验证，仍保留温湿度、视频、烟雾和门禁的完整研究；跨阶段配对数据不足时，仅取消运输先验扩展实验，不删除仓储主线。实验平台、端侧设备和现场合作情况将在研究基础信息补充后据实核验。
